\section{Finite aggregates instead of premature witness selection}
\label{sec:aggregates}

\subsection{A representation with finite, decidable observations}

Work first in $\Q^d$ with the sup metric
$d(x,y)=\max_{i<d}|x_i-y_i|$. This metric is rational-valued and PR on
finite vector codes. Everything below also applies to any fixed coded
rational metric with proved triangle and symmetry laws. The complex
root application uses squared Euclidean comparisons; its estimates
imply the sup-metric estimates because $\|v\|_\infty\le|v|_2$.

For nonempty finite lists $S,T$, repetitions allowed, put
\[
 h(S,T)=\max_{x\in S}\min_{y\in T}d(x,y),\qquad
 H(S,T)=\max\bigl(h(S,T),h(T,S)\bigr).
\]
Minima and maxima mean finite scans with the first index breaking a tie.
All comparisons are rational. Equivalently, $H(S,T)\le r$ says
\[
 (\forall x\in S\ \exists y\in T\ d(x,y)\le r)
 \ \wedge\
 (\forall y\in T\ \exists x\in S\ d(x,y)\le r).
 \tag{H}\label{eq:finite-cover}
\]
This is a bounded formula, hence has a Boolean checker by
Lemma~\ref{lem:booleanization}. Its existential quantifiers are finite
searches, not choices of infinite point names.

\begin{lemma}[Finite aggregate calculus]\label{lem:aggregate-calculus}
The following statements have finite proofs on the supplied lists.
\begin{enumerate}
\item $H(S,U)\le H(S,T)+H(T,U)$.
\item $H(S\mathbin{\|}S',T\mathbin{\|}T')\le
 \max(H(S,T),H(S',T'))$, where $\|$ concatenates lists.
\item With the sup product metric,
 $H(S\times S',T\times T')\le\max(H(S,T),H(S',T'))$.
\item If $d'(f(x),f(y))\le Ld(x,y)$ on the relevant finite inputs,
 then $H(f[S],f[T])\le LH(S,T)$.
\item Writing $d(z,S)=\min_{x\in S}d(z,x)$, one has
 $|d(z,S)-d(w,T)|\le d(z,w)+H(S,T)$.
\end{enumerate}
\end{lemma}

\begin{proof}
For the first assertion, choose for each supplied $x\in S$ a nearest
$y\in T$, then a nearest $z\in U$. The triangle inequality gives
$d(x,z)\le H(S,T)+H(T,U)$. The reverse direction is identical with the
lists reversed. These choices are least-index finite scans.

For concatenation, inspect which part of the concatenated list contains
the current index and use that part's cover. For products, choose the
nearest candidate in each factor and use the sup product metric. For
images, apply the assumed finite Lipschitz inequality to every selected
cover pair. Finally choose a minimizing $y\in T$ for $d(w,T)$ and an
$x\in S$ within $H(S,T)$ of it. Then
$d(z,S)\le d(z,w)+d(w,T)+H(S,T)$. Interchanging the pairs gives the
absolute-value bound. All statements about all list elements are bounded
conjunctions; prefix induction verifies their construction.
\end{proof}

\begin{definition}[Regular aggregate name]\label{def:aggregate-name}
A regular aggregate name is a fixed PR family $k\mapsto S_k$ of
nonempty finite rational lists satisfying
\[
                      H(S_s,S_t)\le e_s+e_t,
             \qquad e_j=2^{-j}.                 \tag{A}\label{eq:aggregate-name}
\]
The family can depend on numerical parameters. Its represented object
is the completion of these finite observations, not a decision procedure
for membership in a finite list of exact points.
\end{definition}

For a fixed rational point $z$, Lemma~\ref{lem:aggregate-calculus}(5)
shows that $k\mapsto d(z,S_k)$ is a regular rational name. Denote its
represented value by $D_S(z)$. If $z$ is itself a regular point name,
use $d(z_{k+1},S_{k+1})$: its change between indices $k,\ell$ is at
most $e_k+e_\ell$. This defines the same distance extension. In semantic
notation,
\[
 |D_S(z)-d(z,S_s)|\le e_s,\qquad
 |D_S(z)-D_S(w)|\le d(z,w).                       \tag{D}\label{eq:distance-name}
\]
The first inequality follows by comparing to $S_t$ and taking an
arbitrarily fine \emph{explicit} later observation. For a finite desired
slack $2^{-k}$, it is enough to take $t\ge k+1$ and a point precision
at least $k+1$. Limits here are shorthand for those finite inequalities.

One useful equality convention for aggregate names is
\[
                     S\simeq T\quad\Longleftrightarrow\quad
              \forall k\ H(S_k,T_k)\le2e_k.     \tag{E}\label{eq:aggregate-equality}
\]
It is an equivalence relation. Only transitivity needs explanation.
If $S\simeq T$ and $T\simeq U$, then for every $v\ge k$,
\[
 H(S_k,U_k)\le2e_k+6e_v.
\]
The left side and $2e_k$ are rationals. If their difference were positive,
its numerator and denominator give an explicit binary-length bound for
$v$ with $6e_v$ smaller than that difference, a contradiction. Thus
$H(S_k,U_k)\le2e_k$. This uses rational decidability, not Markov's
principle. Equivalent names give identical distance names by the same
finite comparison. Equality \eqref{eq:aggregate-equality} is itself
unbounded and is not a Boolean test on program indices.

\subsection{Realizing a point by a prescribed chain}

A fixed PR sequence of finite lists gives more than an arbitrary
qualitative assertion of approximability. It permits the following
explicit point constructor.

\begin{theorem}[Least-neighbor chain]\label{thm:aggregate-chain}
Given a regular aggregate name, an index $s$, and an entry $q\in S_s$,
one can construct a PR point name $x^{s,q}$ such that
\[
                  D_S(x^{s,q})=0,
                  \qquad d(x^{s,q},q)\le3\,2^{-s}.
\]
The construction is uniform in $s$, the entry index, and the numerical
parameters of the fixed list evaluator. It does not assert extensional
or continuous selection of a preferred point from the aggregate.
\end{theorem}

\begin{proof}
Put $q_s=q$. Once $q_j\in S_j$ has been produced, let $q_{j+1}$ be
the first nearest entry of $S_{j+1}$. Nonemptiness and
\eqref{eq:aggregate-name} give
\[
 d(q_j,q_{j+1})\le e_j+e_{j+1}.
\]
For $s\le j<v$, a finite telescoping sum gives
\[
 d(q_j,q_v)\le\sum_{r=j}^{v-1}(e_r+e_{r+1})
            =3e_j(1-2^{-(v-j)})\le3e_j.          \tag{C}\label{eq:chain-tail}
\]
This is induction on a finite table of rational points. Its invariant is
membership of the last point in the last finite list, together with the
displayed sum bound; it does not assert the existence of a limiting point.

Define the output at precision $t$ to be $x^{s,q}(t)=q_{s+t+2}$.
Equation~\eqref{eq:chain-tail} implies, for $u\ge t$,
\[
 d(x^{s,q}(t),x^{s,q}(u))\le3\,2^{-s-t-2}\le2^{-t}.
\]
Thus the output is a regular PR point name. The same finite bound gives
its distance from $q_s$ at most $3e_s$. Membership has a particularly
direct finite verification. Its distance-name sample at index $k$ is
\[
 d\bigl(x^{s,q}(k+1),S_{k+1}\bigr)
 =d(q_{s+k+3},S_{k+1})
 \le e_{s+k+3}+e_{k+1}\le\tfrac58 e_k,
\]
because $q_{s+k+3}\in S_{s+k+3}$. These nonnegative rational samples
therefore name zero. The largest list index needed in this check is
$s+k+3$; all evaluation indices and loop lengths are explicit.

Every step calls the same fixed PR list evaluator and performs a finite
minimum scan. Iteration is primitive recursion on $j$. If an invalid
input produces an empty list, choose the zero vector as a default; the
proof above concerns the stated valid-input hypotheses.
\end{proof}

Define membership in the represented support by $D_S(x)=0$. The theorem
shows that each finite approximant lies within $3e_s$ of a represented
member. Conversely, \eqref{eq:distance-name} shows that every member has
distance at most $e_s$ from $S_s$. This is a statement about its minimum
distance name, not a selection of an exact-radius witness for an arbitrary
real point. Replacing $S_s$ by $S_{s+2}$ absorbs the factor
three if a unit-error convention is desired. This semantic explanation
uses a prescribed PR sequence; it is not an invocation of countable
choice for arbitrary abstract approximation data. Richman's completed
finite-set approach is a useful broader antecedent \cite{richman2000}.

\subsection{Images and optimization, with their exact output meaning}

Finite union and product of regular names remain regular by
Lemma~\ref{lem:aggregate-calculus}. For an approximate represented map,
suppose a fixed PR evaluator $f(q,p)$ obeys the numerical contract
\[
 d'\bigl(f(q,p),f(r,p')\bigr)
          \le Ld(q,r)+2^{-p}+2^{-p'}              \tag{L}\label{eq:approx-lipschitz}
\]
on all rational inputs, with a given nonnegative rational $L$.
A specified invariant domain containing every sample and intermediate
point suffices instead if its membership is included in the contracts.
Choose a natural number $c$ with $2^c\ge L+1$, using integer bit lengths
after clearing the denominator, and put
\[
                  T_k=\bigl[f(q,k+c):q\in S_{k+c}\bigr].
\]
The finite covers and \eqref{eq:approx-lipschitz} give
\[
 H(T_s,T_t)\le(L+1)(e_{s+c}+e_{t+c})\le e_s+e_t.
\]
This proves a PR constructor for the completed finite-image object.
The finite observation indices are exactly the displayed ones. If $S\simeq S'$, the same calculation gives $H(T_k,T'_k)\le2(L+1)e_{k+c}\le2e_k$, so the finite-image constructor is extensional on aggregate names.

The action on represented point names is also explicit:
\[
             \widehat f(x)(k)=f(x(k+c),k+c).
\]
Contract \eqref{eq:approx-lipschitz} and regularity of $x$ bound its
sample difference by $(L+1)(e_{s+c}+e_{t+c})\le e_s+e_t$.
Equivalent point names give equivalent outputs by the same estimate,
and $d'(\widehat f(x),\widehat f(y))\le Ld(x,y)$ in represented
notation. Thus later occurrences of $f(x)$ refer to this prescribed
extension, rather than to a separately assumed completed function.

The word ``image'' does not conceal a preimage selection principle.
If $D_S(x)=0$, then $D_T(f(x))=0$ by the Lipschitz contract.
We have not proved a uniform constructor of $x$ from an arbitrary named
$y$ with $D_T(y)=0$. Likewise, the distance for a union is
$\min(D_S,D_T)$; a proof that it vanishes is not a decision of the
disjunction $D_S=0\vee D_T=0$. This distinction lets aggregate
representations preserve useful information without reintroducing the
selection difficulty they were meant to avoid.

\begin{theorem}[Value and approximate witnesses for optimization]
\label{thm:aggregate-optimization}
Let $f$ be a real-valued map satisfying \eqref{eq:approx-lipschitz}.
Its infimum over the represented support has a PR regular name.
For each $k$, one can produce a PR point name $x^{(k)}$ in that support
whose value is within $2^{-k}$ above the infimum. No one exact minimizing
point is selected by this theorem. The same assertions hold for the
supremum by replacing $f$ with $-f$.
\end{theorem}

\begin{proof}
Let $a_s=\min_{q\in S_s}f(q,s)$. Pick finite minimizers and use the
two directed covers. Contract \eqref{eq:approx-lipschitz} gives
\[
                         |a_s-a_t|\le(L+1)(e_s+e_t).
\]
Thus $a_{k+c}$, with the same $c$ as above, is a regular PR name;
write its value as $a$. In particular $|a-a_s|\le(L+1)e_s$.
If $x$ belongs to the support, choose a nearest finite point to a
sufficiently accurate rational observation of $x$. Equations
\eqref{eq:distance-name} and \eqref{eq:approx-lipschitz} imply
$a\le f(x)$ after making their explicit errors arbitrarily small.
For a desired slack $\varepsilon$, it suffices to choose $s$ with
$(4L+2)e_s<\varepsilon$ and point precision at least $s$.
This proves the lower-bound part directly from finite observations.

Choose $q_s\in S_s$ with $f(q_s,s)=a_s$ by a finite scan. The chain
theorem gives $x_s$ in the support with $d(x_s,q_s)\le3e_s$.
Writing $f(q_s)$ for its represented value, the same contract yields
\[
 f(x_s)\le f(q_s)+3Le_s
          \le a_s+(3L+1)e_s
          \le a+(4L+2)e_s.
\]
Choose $b$ with $2^b\ge4L+2$ and use $s=k+b$. Then $x^{(k)}=x_s$
has the required error. Both the finite minimizing index and the
subsequent point-name evaluator are PR. The approximating minimizers can
switch components as $k$ changes; nothing in this proof asserts their
coherence as a single minimizing name.
\end{proof}

\subsection{A reusable application to spectra}

The root-support producer in Section~\ref{sec:rootsets} combines with
this calculus to give an aggregate spectrum for an $n$ by $n$ matrix
of PR complex names, $n\ge1$. This statement concerns the characteristic
polynomial and its root support; it does not require eigenvectors,
rank decisions, or a Jordan normal form.

Here is an explicit coefficient precision budget. Use complex entry
names regular in the Euclidean norm. Suppose every sampled
entry has modulus at most an integer $B\ge1$; one can take
$B=3+\lceil\max_{i,j}\|a_{ij}(0)\|_1\rceil$.
For names given with a sup-coordinate convention, sample one extra
precision bit to obtain the Euclidean convention. A characteristic
coefficient is a signed sum of principal minors. There are at most
$2^n$ such minors in any coefficient and at most $n!$ determinant terms
per minor. A product of $r\le n$ entries changes by at most
$rB^{r-1}\delta$ when corresponding entries change by at most $\delta$;
expand the product difference one factor at a time. Thus the safe bound
\[
 A=2^n n!(B+1)^n,\qquad
 C=2^n n!\,n(B+1)^{n-1}
\]
respectively bounds every lower coefficient and its sensitivity to the
maximum entry error. Compute the characteristic coefficients by finite
determinant sums, using entry precision $s+c$ with $2^c\ge C$.
Their Cauchy error is then at most $e_s+e_t$. All loops and factorials
are PR uniformly in $n$.

Apply the root theorem, then the maximum construction to $z\mapsto|z|$.
Rational square-root bisection supplies approximations to $|z|$, with
the inequality $||z|-|w||\le|z-w|_2\le2\|z-w\|_\infty$;
thus use the safe Lipschitz constant $L=2$ in the sup-metric calculus.
The relevant disk bound is explicit.
This gives a PR name for the spectral radius, and finite coverage of
every spectral point, including at repeated eigenvalues. The result is
a direct application of our proved interfaces, not a claim to have
formalized the spectral theorem. The coefficient calculations are
compatible with the generic-matrix approach of the Cayley--Hamilton
packet \cite{repo-cayley}.

\subsection{Lean realization of the finite interface}

\path{lean/FiniteHausdorff.lean} defines directed covers, a finite
Boolean checker, the two-sided relation, and regular aggregate names.
It proves checker reflection, relational cover composition, finite union,
finite image transport, and the rational-metric triangle theorem.
Its exact axiom audit is included: the purely relational composition
lemma uses no axioms, while the rational triangle instance inherits
\texttt{Classical.choice} from Lean's standard rational-order lemmas.
This is a native-library dependency, not an added object-iPRA axiom.
The human finite rational proof above and its future arithmetization
remain independent obligations.

The companion \path{lean/AggregateTargets.lean} supplies concrete
rational-vector distance, finite minimum, least-neighbor iteration,
distance-name membership, represented-map evaluation, and approximate
minimizer programs. Their full chain, value, and witness theorems are
named unproved propositions. These definitions specify the remaining
mathematical statements rather than assuming them as structure fields.

To finish the object proof, encode the metric and lists, establish the
arithmetic triangle law by rational cross multiplication, prove the
finite-minimum specification by a prefix invariant, and replace every
list quantifier in \eqref{eq:finite-cover} by its Boolean fold. For the
chain theorem, the invariant is exactly membership plus
\eqref{eq:chain-tail}. For optimization, use the finite minimum and
the displayed $(L+1)$ and $(4L+2)$ budgets. No compactness theorem from
an ambient classical real analysis library is needed.
